\documentclass[a4paper,11pt]{article}

\usepackage[utf8]{inputenc}
\usepackage[T2A]{fontenc}
\usepackage[english,russian]{babel}
\usepackage[margin=2.5cm]{geometry}
\usepackage{amsmath,amssymb,amsthm,mathtools}
\usepackage{enumitem}
\usepackage{booktabs}
\usepackage{xcolor}
\usepackage[colorlinks=true,linkcolor=blue,urlcolor=blue,citecolor=blue]{hyperref}
\usepackage{url}
\newcommand{\lean}[1]{\mbox{\path{#1}}}

\theoremstyle{plain}
\newtheorem{theorem}{Теорема}[section]
\newtheorem{lemma}[theorem]{Лемма}
\newtheorem{proposition}[theorem]{Утверждение}
\newtheorem{corollary}[theorem]{Следствие}
\theoremstyle{definition}
\newtheorem{definition}[theorem]{Определение}
\theoremstyle{remark}
\newtheorem*{remark}{Замечание}

\newcommand{\R}{\mathbb{R}}
\newcommand{\Q}{\mathbb{Q}}
\newcommand{\abs}[1]{\left\lvert #1 \right\rvert}
\newcommand{\set}[1]{\left\{ #1 \right\}}
\newcommand{\ip}[2]{\langle #1, #2\rangle}
\newcommand{\Rpp}[1][2]{\R^{#1}_{++}}
\newcommand{\Rp}[1][2]{\R^{#1}_{+}}
\DeclareMathOperator{\conv}{conv}
\DeclareMathOperator{\cone}{cone}
\DeclareMathOperator{\Sp}{Sp}
\DeclareMathOperator{\Out}{Out}
\DeclareMathOperator{\Bnd}{Bnd}
\newcommand{\WS}{\textsc{WS}}
\newcommand{\COL}{\textsc{Col}}
\newcommand{\Spm}{\Sp_\star}
\newcommand{\gs}{\mathrel{\succ}}
\newcommand{\ngs}{\mathrel{\not\succ}}
\newcommand{\e}{\mathbf e}

\title{Разрешимость на плоскости за полиномиальное время\\ и точный алгоритм минимального возмущения цен}
\author{}
\date{11 октября 2026}

\begin{document}
\maketitle

\begin{abstract}
При $d=2$ проверка неоклассической и слабой разрешимости при фиксированных ценах сводится к одной задаче линейного программирования: потоку в ациклическом графе с $O(T^3)$ дугами. Доказательство полностью проверено в Lean~4, включая разложение потока на пути. Отсюда точная и полиномиальная процедура разделения для задачи о минимальном возмущении цен. На ней построен точный алгоритм отсечений для любой выпуклой сепарабельной стоимости: $\abs{\log(q/p)}$, возмущения $\pi$ из заметок, $\abs{q/p-1}$. Нижняя граница~--- значение мастер-задачи, решённой методом ветвей и границ без внешнего решателя. Верхняя граница проверяется в точной рациональной арифметике. На шести плоских рядах Berndt--Wood ($T=25$) алгоритм находит $\rho$ с зазором порядка $10^{-6}$ для всех трёх стоимостей. Для $\pi$ значения совпадают с прежними сертифицированными.
\end{abstract}

\section{Постановка}

Обозначения взяты из заметок (\path{latex/notes.tex}) и работы \path{complexity/complexity.tex}. Цены $P=\set{p_t\in\Rpp}_{t\in[T]}$, $X_t:=p^1_t$, $Y_t:=p^2_t$. Рассматриваются два \emph{режима} $\star$: слабый и строгий. Запись $u\gs0$ означает $u\ge0$ в слабом режиме и $u>0$ в строгом. Множество $S\subset[T]$ лежит в $\Spm(P)$, если для каждого $t\in S$ существует ненулевой $\xi\in\Rp$ с $\ip{\xi}{p_s-p_t}\gs0$ при всех $s\notin S$. Так $\Spm(P)$~--- это $\widehat{\Sp}(P)$ в слабом режиме и $\Sp(P)$ в строгом. Вектор $y\in\Rp[T]$ \emph{разрешим в режиме $\star$}, если $y\in\cone\set{\mathbf 1_S\mid S\in\Spm(P)}$. По теореме о разрешимости и критерию слабой разрешимости из заметок это неоклассическая и слабая разрешимость.

Работа состоит из двух частей.
\begin{itemize}
    \item Задача $\WS_\star(2)$: по рациональным $P\subset\Rpp$ и $y$ решить, разрешим ли $y$ в режиме $\star$. Раздел~\ref{sec:lp}: она решается одной задачей линейного программирования полиномиального размера.
    \item Задача о минимальном возмущении (раздел~\ref{sec:pert}): найти ближайшие к $P$ цены $Q$, при которых $y$ слабо разрешим. Расстояние задаётся выпуклой сепарабельной стоимостью.
\end{itemize}

\section{Цепочки и полосы}\label{sec:chains}

Положим
\[
    \chi(a,b,q):=(X_b-X_a)(Y_q-Y_a)-(Y_b-Y_a)(X_q-X_a).
\]
При $X_a<X_b$ число $\chi(a,b,q)$ положительно, когда $p_q$ лежит строго выше прямой через $p_a$ и $p_b$.

\begin{definition}[цепочка]\label{def:chain}
    \emph{Цепочка}~--- последовательность индексов $v_0,\dots,v_k$ с $X_{v_0}<\dots<X_{v_k}$, $Y_{v_0}>\dots>Y_{v_k}$ и $\chi(v_j,v_{j+1},v_{j+2})>0$ при $j+2\le k$. Индекс $q$ лежит в \emph{области} $C(V)$, если $X_q\ge X_{v_0}$, $Y_q\ge Y_{v_k}$ и $\chi(v_j,v_{j+1},q)\ge0$ при всех $j<k$, и во \emph{внутренности} $C^\circ(V)$, если все эти неравенства строгие. \emph{Запрещённая область} $F_\star(V)$~--- это $C^\circ(V)$ в слабом режиме и $C(V)$ в строгом.
\end{definition}

Область $C(V)$~--- это $\conv\set{p_v}_{v\in V}+\Rp$.

\begin{theorem}[цепочки, \cite{Complexity}]\label{thm:chain}
    $S\in\Spm(P)$ тогда и только тогда, когда $S=[T]$ или существует цепочка $V$, для которой вершины $V$ не лежат в $S$, каждый индекс вне $C(V)$ лежит в $S$, а каждый индекс из $F_\star(V)$ не лежит в $S$.
\end{theorem}

Доказательство дано в~\cite{Complexity} и проверено в Lean (\lean{isReachableB_iff_chain}). Кратко: свидетель индекса $t\in S$~--- это $(1,0)$, $(0,1)$ или нормаль ребра $(Y_{v_j}-Y_{v_{j+1}},\,X_{v_{j+1}}-X_{v_j})$, смотря какое неравенство области нарушает $p_t$. Обратно, цепочку для $U=[T]\setminus S$ строит «заворачивание подарка», а свидетель точки запрещённой области противоречит неравенствам в вершинах.

Индексы вне $V$ делятся на \emph{полосы}: левую $X_q\le X_{v_0}$, полосу ребра $j$ $X_{v_j}<X_q\le X_{v_{j+1}}$ и правую $X_q>X_{v_k}$. Каждый индекс вне $V$ лежит ровно в одной полосе.

\begin{lemma}[локальность, \cite{Complexity}]\label{lem:local}
    В полосе ребра $j$: $q\in C(V)\iff\chi(v_j,v_{j+1},q)\ge0$ и $q\in C^\circ(V)\iff\chi(v_j,v_{j+1},q)>0$. В левой полосе $q\notin C^\circ(V)$, и $q\in C(V)\iff X_q=X_{v_0}$, $Y_q\ge Y_{v_0}$. В правой полосе $q\in C(V)\iff Y_q\ge Y_{v_k}$ и $q\in C^\circ(V)\iff Y_q>Y_{v_k}$.
\end{lemma}

Строгая выпуклость нужна именно здесь: принадлежность области решается одним неравенством полосы.

\section{Разрешимость как задача о потоке}\label{sec:lp}

\begin{definition}[граф цепочек]\label{def:graph}
    Вершины графа $G_P$: источник $\sigma$, сток $\tau$, вершины $A_b$ при $b\in[T]$ («цепочка начинается в $b$») и $E_{ab}$ при $X_a<X_b$, $Y_a>Y_b$ («последнее ребро цепочки~--- $ab$»). Дуги:
    \[
        \sigma\to\tau,\quad \sigma\to A_b,\quad A_a\to E_{ab},\quad E_{ab}\to E_{bc}\ \text{при}\ X_b<X_c,\ Y_b>Y_c,\ \chi(a,b,c)>0,\quad A_b\to\tau,\quad E_{ab}\to\tau.
    \]
    Каждая дуга $e$ несёт два множества индексов: $\Out_e$ (вне области) и $\Bnd_e$ (на границе).
    \begin{itemize}
        \item Дуга в $A_b$ (левая полоса $b$): $\Out_e=\set{q\ne b\mid X_q<X_b\ \text{или}\ X_q=X_b,\,Y_q<Y_b}$, $\Bnd_e=\set{q\ne b\mid X_q=X_b,\,Y_q\ge Y_b}$.
        \item Дуга в $E_{ab}$ (полоса ребра $ab$): $\Out_e=\set{q\ne b\mid X_a<X_q\le X_b,\ \chi(a,b,q)<0}$, $\Bnd_e$~--- то же с $\chi(a,b,q)=0$.
        \item Дуга из $A_b$ или $E_{ab}$ в $\tau$ (правая полоса $b$): $\Out_e=\set{q\mid X_q>X_b,\ Y_q<Y_b}$, $\Bnd_e=\set{q\mid X_q>X_b,\ Y_q=Y_b}$.
        \item Дуга $\sigma\to\tau$: $\Out_e=[T]$, $\Bnd_e=\varnothing$.
    \end{itemize}
    Для пути $\pi$ из $\sigma$ в $\tau$ положим $\Out(\pi):=\bigcup_{e\in\pi}\Out_e$ и $\Bnd(\pi):=\bigcup_{e\in\pi}\Bnd_e$.
\end{definition}

Граф ацикличен: вдоль дуг $E_{ab}\to E_{bc}$ координата $X$ последней вершины растёт. Путь $\sigma\to A_{v_0}\to E_{v_0v_1}\to\dots\to E_{v_{k-1}v_k}\to\tau$ соответствует цепочке $v_0,\dots,v_k$, и каждый путь, кроме $\sigma\to\tau$, имеет такой вид. Вершин $O(T^2)$, дуг $O(T^3)$. Положим $\beta_\star:=1$ в слабом режиме и $\beta_\star:=0$ в строгом.

\begin{theorem}[семейство]\label{thm:family}
    $S\in\Spm(P)$ тогда и только тогда, когда для некоторого пути $\pi$ из $\sigma$ в $\tau$
    \[
        \Out(\pi)\subset S\subset\Out(\pi)\cup\beta_\star\Bnd(\pi).
    \]
    Здесь $\beta_\star\Bnd(\pi)$ означает $\Bnd(\pi)$ при $\beta_\star=1$ и $\varnothing$ при $\beta_\star=0$.
\end{theorem}

\begin{proof}
    Путь $\sigma\to\tau$ даёт $\Out=[T]$, то есть $S=[T]$. Для пути цепочки $V$ по лемме~\ref{lem:local} и разбиению на полосы
    \[
        \Out(\pi)=\set{q\mid q\notin C(V)},\qquad \Bnd(\pi)=\set{q\notin V\mid q\in C(V)\setminus C^\circ(V)}.
    \]
    Вершины лежат в $C(V)$, но не в $C^\circ(V)$. Поэтому три условия теоремы~\ref{thm:chain} равносильны включениям. В слабом режиме запрещена внутренность, и $S\subset\Out(\pi)\cup\Bnd(\pi)$ означает, что в $S$ нет вершин и точек внутренности. В строгом режиме запрещена вся область, и $S=\Out(\pi)$.
\end{proof}

\begin{lemma}[полосы не пересекаются]\label{lem:disjoint}
    Вдоль пути множества $\Out_e$ и $\Bnd_e$ разных дуг попарно не пересекаются, и $\Out(\pi)\cap\Bnd(\pi)=\varnothing$. Поэтому $\mathbf 1_{\Out(\pi)}=\sum_{e\in\pi}\mathbf 1_{\Out_e}$ и $\mathbf 1_{\Bnd(\pi)}=\sum_{e\in\pi}\mathbf 1_{\Bnd_e}$.
\end{lemma}

\begin{proof}
    Множества дуги лежат в её полосе, полосы разных дуг пути не пересекаются. Внутри одной полосы $\Out_e$ и $\Bnd_e$ различаются знаком $\chi$ или сравнением $Y$.
\end{proof}

\begin{lemma}[разложение потока]\label{lem:decomp}
    Пусть в конечном графе есть ранг, строго растущий вдоль дуг, в источник не входят дуги и из стока не выходят. Тогда всякий поток (неотрицательный, сохраняющийся во всех вершинах, кроме источника и стока) есть неотрицательная комбинация путей из источника в сток.
\end{lemma}

\begin{proof}
    Индукция по числу дуг с ненулевым потоком. Если поток ненулевой, то из источника выходит дуга с положительным потоком. Действительно, возьмём положительную дугу $u\to v$ с наименьшим рангом $u$. Если $u$~--- не источник, то $u$~--- и не сток, и по сохранению в $u$ входит положительная дуга с меньшим рангом начала. Далее идём по положительным дугам. Ранг растёт, поэтому путь конечен, а сохранение не даёт остановиться раньше стока. Каждая дуга пути встречается в нём один раз. Вычтем путь с весом $m$, равным наименьшему потоку на его дугах. Остаток~--- поток, так как путь сохраняется во внутренних вершинах, и на дуге, где достигнут минимум, поток стал нулевым.
\end{proof}

\begin{lemma}[куб]\label{lem:cube}
    Пусть $O\cap B=\varnothing$ и $0\le z\le\mathbf 1_B$. Тогда $\mathbf 1_O+z\in\cone\set{\mathbf 1_{O\cup B'}\mid B'\subset B}$.
\end{lemma}

\begin{proof}
    Индукция по $B$. Для $B=B_0\cup\set q$ запишем $\mathbf 1_O+z=(1-z_q)(\mathbf 1_O+z|_{B_0})+z_q(\mathbf 1_{O\cup\set q}+z|_{B_0})$ и применим предположение к $O$ и к $O\cup\set q$.
\end{proof}

\begin{theorem}[$d=2$]\label{thm:lp}
    Вектор $y$ разрешим в режиме $\star$ тогда и только тогда, когда разрешима система
    \[
        f\ \text{--- поток из}\ \sigma\ \text{в}\ \tau\ \text{в}\ G_P,\qquad y=\sum_ef_e\mathbf 1_{\Out_e}+g,\qquad 0\le g\le\beta_\star\sum_ef_e\mathbf 1_{\Bnd_e}.
    \]
\end{theorem}

\begin{proof}
    $(\Rightarrow)$ Пусть $y=\sum_Sm_S\mathbf 1_S$, $m_S>0$, $S\in\Spm(P)$. По теореме~\ref{thm:family} для каждого $S$ выберем путь $\pi_S$. Положим $f:=\sum_Sm_S\mathbf 1_{\pi_S}$ и $g:=\sum_Sm_S\mathbf 1_{S\setminus\Out(\pi_S)}$. Сумма путей~--- поток. По лемме~\ref{lem:disjoint} $\mathbf 1_S=\sum_{e\in\pi_S}\mathbf 1_{\Out_e}+\mathbf 1_{S\setminus\Out(\pi_S)}$. В слабом режиме $S\setminus\Out(\pi_S)\subset\Bnd(\pi_S)$, откуда $g\le\sum_ef_e\mathbf 1_{\Bnd_e}$. В строгом $S=\Out(\pi_S)$ и $g=0$.

    $(\Leftarrow)$ По лемме~\ref{lem:decomp} $f=\sum_im_i\mathbf 1_{\pi_i}$, и по лемме~\ref{lem:disjoint} $\sum_ef_e\mathbf 1_{\Out_e}=\sum_im_i\mathbf 1_{\Out(\pi_i)}$, $h:=\sum_ef_e\mathbf 1_{\Bnd_e}=\sum_im_i\mathbf 1_{\Bnd(\pi_i)}$. Разделим $g$ пропорционально: $z_i(q):=g_q/h_q$ при $q\in\Bnd(\pi_i)$ и $0$ иначе. Если $h_q=0$, то $g_q=0$. Тогда $0\le z_i\le\beta_\star\mathbf 1_{\Bnd(\pi_i)}$ и $\sum_im_iz_i=g$, так что $y=\sum_im_i(\mathbf 1_{\Out(\pi_i)}+z_i)$. По лемме~\ref{lem:cube} каждое слагаемое лежит в конусе векторов $\mathbf 1_{\Out(\pi_i)\cup B'}$ с $B'\subset\beta_\star\Bnd(\pi_i)$, а эти множества лежат в $\Spm(P)$ по теореме~\ref{thm:family}.
\end{proof}

\begin{corollary}[полиномиальность]\label{cor:poly}
    В любом режиме $\WS_\star(2)$ решается за полиномиальное время. Граф строится за $O(T^4)$ арифметических операций над числами вида $\chi(a,b,q)$, степень которых по входным числам равна $2$. Система теоремы~\ref{thm:lp} имеет $O(T^3)$ переменных, $O(T^2)$ ограничений, коэффициенты $0$, $\pm1$ и правую часть $y$, поэтому решается за полиномиальное время методом эллипсоидов или внутренних точек. По базисному решению за полиномиальное время находится и представление $y=\sum_km_k\mathbf 1_{S_k}$ с $S_k\in\Spm(P)$: разложение потока даёт не больше числа дуг путей, лемма~\ref{lem:cube} и теорема Каратеодори сокращают их до $T$ множеств.
\end{corollary}

\begin{corollary}[задача о столбце]\label{cor:col}
    $\max\set{\sum_{t\in S}w_t\mid S\in\Spm(P)}$ равен длиннейшему пути в $G_P$ с весами дуг $\sum_{q\in\Out_e}w_q+\beta_\star\sum_{q\in\Bnd_e}\max(w_q,0)$. Время $O(T^4)$.
\end{corollary}

\begin{proof}
    По теореме~\ref{thm:family} при фиксированном пути лучшее $S$ берёт $\Out(\pi)$ и граничные точки с $w_q>0$. Вес пути аддитивен по лемме~\ref{lem:disjoint}.
\end{proof}

\begin{corollary}[сертификат неразрешимости]\label{cor:farkas}
    Если система теоремы~\ref{thm:lp} неразрешима, двойственная задача даёт $w\in\R^T$ с $\ip{w}{y}<0$ и $\sum_{t\in S}w_t\ge0$ для всех $S\in\Spm(P)$. Его проверка~--- одна задача о столбце.
\end{corollary}

\begin{remark}
    Теорема~\ref{thm:lp} заменяет перебор $2^T$ множеств из алгоритма заметок. Без строгой выпуклости локальность ломается: для невыпуклой ломаной неравенство своей полосы не решает принадлежность области.
\end{remark}

\section{Минимальное возмущение цен на плоскости}\label{sec:pert}

\subsection{Постановка}

Возмущение измеряется сепарабельной стоимостью
\[
    c(P,Q):=\sum_{t\in[T],\,i\in\set{1,2}}\varphi\bigl(r(p^i_t,q^i_t)\bigr),
\]
где $r(p,q)=\log(q/p)$ (\emph{логарифмическая шкала}) или $r(p,q)=q/p$ (\emph{линейная шкала}), а $\varphi$ выпукла и неотрицательна. Примеры:
\begin{itemize}
    \item $\varphi(r)=\abs r$ в логарифмической шкале: $c=\sum\abs{\log(q/p)}$, симметричная метрика;
    \item $\varphi(r)=e^{\abs r}-1$ в логарифмической шкале: $c=\sum\pi(p,q)$, возмущение $\pi(p,q)=\max(q/p,p/q)-1$ из заметок, выпукла и в линейной шкале;
    \item $\varphi(r)=\abs{r-1}$ в линейной шкале: $c=\sum\abs{q/p-1}$, то есть $\sum\abs\varepsilon$ при $q=(1+\varepsilon)p$.
\end{itemize}
Положим
\[
    \rho_c(P,y):=\inf\set{c(P,Q)\mid Q\in\Rpp[2\times T],\ y\ \text{разрешим при ценах}\ Q}.
\]

\begin{proposition}[минимум достигается]\label{prop:min}
    Пусть $\varphi$ непрерывна и $\varphi(r)\to\infty$ при $\abs r\to\infty$ в логарифмической шкале. Тогда
    \[
        \rho_c(P,y)=\min\set{c(P,Q)\mid Q\in\Rpp[2\times T],\ y\ \text{слабо разрешим при ценах}\ Q},
    \]
    и минимум достигается.
\end{proposition}

\begin{proof}
    Цены на гиперболе $q^1q^2=\kappa$ дают слабо разрешимый $y$ (лемма о гиперболе в заметках), так что множество непусто. Подуровень $c\le c_0$ лежит в коробке $p\,e^{-M}\le q\le p\,e^{M}$, а слабо разрешимые цены образуют замкнутое множество (по определению это замыкание). Компактность даёт минимум (Lean \lean{exists_min_cost}). Равенство с точной нижней гранью по разрешимым ценам следует из теоремы о замыкании и непрерывности $c$, как в лемме о достижении минимума заметок.
\end{proof}

Для линейной шкалы стоимость $\abs{q/p-1}$ ограничена при $q\to0$, и минимум может уходить на границу $\Rp[2\times T]$. В вычислениях координаты ограничены снизу: $q\ge10^{-6}p$.

\subsection{Блокировка на плоскости}

Индекс $t$ \emph{заблокирован} множеством $R$ при ценах $Q$, если существует ненулевой $\xi\in\Rp$ с $\ip{\xi}{q_r-q_t}\ge0$ при всех $r\in R$. Иначе $t$ \emph{строго покрыт} множеством $R$. Если $S\in\widehat{\Sp}(Q)$, $t\in S$ и $R\cap S=\varnothing$, то $t$ заблокирован множеством $R$.

\begin{lemma}[Хелли на отрезке]\label{lem:helly}
    $t$ заблокирован множеством $R$ тогда и только тогда, когда $t$ заблокирован каждой парой $\set{r,s}\subset R$, $r=s$ допускается.
\end{lemma}

\begin{proof}
    Свидетеля нормируем: $\xi=(\theta,1-\theta)$, $\theta\in[0,1]$. Каждое $r$ допускает отрезок значений $\theta$, а попарно пересекающиеся отрезки прямой имеют общую точку.
\end{proof}

\begin{lemma}[блокировка парой]\label{lem:pair}
    $t$ заблокирован парой $\set{r,s}$ тогда и только тогда, когда выполнено одно из условий:
    \begin{enumerate}[label=(\alph*)]
        \item $X_t\le X_r$ и $X_t\le X_s$;
        \item $Y_t\le Y_r$ и $Y_t\le Y_s$;
        \item $X_r\le X_s$, $Y_s\le Y_r$ и $q_t\le\mu q_r+(1-\mu)q_s$ при некотором $\mu\in[0,1]$;
        \item то же с переставленными $r$ и $s$.
    \end{enumerate}
\end{lemma}

\begin{proof}
    В случаях (a), (b) годятся $\xi=(1,0)$ и $(0,1)$. В случае (c) годится нормаль отрезка $\xi=(Y_r-Y_s,\,X_s-X_r)\ge0$: $\ip{\xi}{q_r}=\ip{\xi}{q_s}\ge\ip{\xi}{q_t}$. Обратно, пусть $a:=q_r-q_t$, $b:=q_s-q_t$ и свидетель $\xi$. Если (a) и (b) нарушены, знаки координат $a$ и $b$ перекрёстные, скажем $a_1<0\le a_2$ и $b_2<0\le b_1$. Тогда $\xi_1>0$, и из $\ip{\xi}{a}\ge0$, $\ip{\xi}{b}\ge0$ следует $a_1b_2-a_2b_1\ge0$, то есть $q_t$ не выше прямой $q_rq_s$. При $\mu:=(X_s-X_t)/(X_s-X_r)$ получаем (c).
\end{proof}

\begin{lemma}[монотонность]\label{lem:mono}
    Если $t$ заблокирован множеством $R$ при ценах $Q$, $q'_t\le q_t$ и $q'_r\ge q_r$ при $r\in R$, то $t$ заблокирован множеством $R$ при ценах $Q'$.
\end{lemma}

\begin{corollary}[покрытие парой]\label{cor:cover}
    Если $S\notin\widehat{\Sp}(Q)$, то некоторое $t\in S$ строго покрыто множеством $R\subset[T]\setminus S$ из не более чем двух индексов.
\end{corollary}

\subsection{Отсечения}

\emph{Отсечение}~--- конечный набор $C$ пар $(t,R)$ с $\abs R\le2$. Цены $Q$ \emph{удовлетворяют} отсечению, если хотя бы для одной пары $t$ заблокирован множеством $R$.

\begin{lemma}[отсечение, \cite{Pert}]\label{lem:cut}
    Пусть $\ip{w}{y}<0$ и всякое $S$ с $\sum_{t\in S}w_t<0$ содержит $t$ и не пересекает $R$ для некоторой пары $(t,R)\in C$. Тогда все цены, при которых $y$ слабо разрешим, удовлетворяют отсечению $C$.
\end{lemma}

Утверждение проверено в Lean (\lean{exists_isBlocked_of_cut}). Важно, что условие на $C$ комбинаторное и не зависит от цен.

\paragraph{Построение отсечения.} Пусть $y$ не слабо разрешим при ценах $Q^*$. Следствие~\ref{cor:farkas} даёт $w$. Отсечение строится жадно. Пока есть $S$ с $\sum_{t\in S}w_t<0$, не задетое парами из $C$ (это задача с булевыми переменными $x\in\set{0,1}^T$: $\min\ip{w}{x}$ при $x_t\le\sum_{r\in R}x_r$ для $(t,R)\in C$), добавляем в $C$ строгое покрытие $(t,R)$ множества $S$ из следствия~\ref{cor:cover}. Множество $S$ с $\sum_Sw_t<0$ не лежит в $\widehat{\Sp}(Q^*)$, поэтому такое покрытие есть. Все пары $C$ строго покрывают при $Q^*$: \emph{новое отсечение нарушено в $Q^*$}.

\subsection{Мастер-задача и метод ветвей и границ}

Для набора отсечений $\mathcal C$ \emph{мастер-задача}
\[
    M(\mathcal C):=\min\set{c(P,Q)\mid Q\ \text{удовлетворяет всем}\ C\in\mathcal C}.
\]
По лемме~\ref{lem:cut} и утверждению~\ref{prop:min} $M(\mathcal C)\le\rho_c(P,y)$. Начальные отсечения~--- доминирование: при $y_t>y_s$ отсечение $\set{(t,\set s)}$.

Мастер решается методом ветвей и границ. Узел задаётся
\begin{itemize}
    \item ограничениями порядка $q^i_a\le q^i_b$;
    \item требованиями отрезка $(t,r,s)$: условие (c) леммы~\ref{lem:pair} с параметром $\mu$ в отрезке $[\mu^-,\mu^+]$;
    \item коробкой для координат: $q^{i,-}_t\le q^i_t\le q^{i,+}_t$.
\end{itemize}
\emph{Ветвление по отсечению.} Если решение узла нарушает отсечение, для каждой его пары $(t,R)$ создаются потомки по формам блокировки леммы~\ref{lem:pair}: (a) и (b) добавляют ограничения порядка, (c) и (d) добавляют ориентацию отрезка, следствия $X_t\le X_s$, $Y_t\le Y_r$ (для (c)) и требование отрезка. Для $R=\set r$ потомки~--- $X_t\le X_r$ и $Y_t\le Y_r$. По лемме~\ref{lem:pair} потомки покрывают все допустимые точки узла.

\emph{Релаксация узла}~--- задача линейного программирования в переменных шкалы: $z=\log q$ или $z=q$. Стоимость заменена касательными минорантами $\varphi$. Для кусочно-линейной $\varphi$ это точно, для гладкой касательные добавляются в точках решения, пока зазор не меньше $10^{-7}$. Ограничения порядка линейны в любой шкале. Требование отрезка релаксируется так:
\begin{itemize}
    \item произведения $\mu q^i_r$ и $(1-\mu)q^i_s$~--- оценками сверху Маккормика по коробкам $\mu$ и $q$;
    \item в логарифмической шкале $q=e^z$: для $q_t$~--- касательными к $e^z$ снизу, они тоже добавляются в точках решения; для $q_r$, $q_s$~--- секущими по коробке сверху.
\end{itemize}
Каждая точка узла, удовлетворяющая его требованиям, допустима для релаксации. Поэтому значение релаксации~--- нижняя граница узла. Без требований отрезка релаксация распадается на две задачи изотонной регрессии, по одной на координату. Возмущение $\pi$ выпукло и в линейной шкале. В ней требование отрезка при фиксированном $\mu$ линейно, поэтому для $\pi$ берётся линейная шкала.

\emph{Пространственное ветвление.} Если в решении узла требование отрезка нарушено, а ограничения порядка выполнены, делится коробка $\mu$ или одной из координат $r$, $s$: та, на которую приходится наибольшая ошибка оценки (Маккормика или секущей) в решении. Делится она в значении из решения LP. Тогда это значение становится концом обеих коробок, и оценки там точны. Если значение близко к краю, коробка делится пополам.

\emph{Отсев.} Узел с требованием $(t,r,s)$ отбрасывается, если $t$ не заблокирован парой $\set{r,s}$ в углу коробки $(q_t^-,q_r^+,q_s^+)$. По лемме~\ref{lem:mono} тогда $t$ не заблокирован ни в одной точке коробки. Отсев точный.

\begin{proposition}\label{prop:bb}
    Метод ветвей и границ, выбирающий узел с наименьшей границей и останавливающийся на первом узле, решение которого удовлетворяет всем отсечениям с допуском $\tau$, возвращает число не больше $M(\mathcal C)$. При сжатии коробок релаксация отрезка стремится к точному условию, поэтому при $\tau>0$ метод конечен.
\end{proposition}

\begin{proof}
    Оптимум мастера лежит в одном из листьев. Отсев не теряет допустимых точек (лемма~\ref{lem:mono}), потомки покрывают узел (лемма~\ref{lem:pair}), граница узла не больше стоимости любой его допустимой точки. Узел с наименьшей границей поэтому не дороже оптимума. Ошибки оценок Маккормика, касательных и секущих стремятся к нулю вместе с шириной коробок.
\end{proof}

\subsection{Алгоритм}

\begin{enumerate}
    \item Верхняя граница: цены на гиперболе $q^1_t=p^1_t$, $q^2_t=\kappa/p^1_t$ с лучшим $\kappa$. Коробка мастера: $\varphi(r)\le$ верхней границы.
    \item Решить мастер, получить $Q^*$ и нижнюю границу $M(\mathcal C)$.
    \item Если $y$ слабо разрешим при $Q^*$ (теорема~\ref{thm:lp}), $Q^*$ оптимален. Точная проверка: сдвиг леммы о наклоне из заметок, $q_t\mapsto q_t+\delta\bigl(\sum_ie^{-q^i_t}-\delta y_t\bigr)\e$ с малым $\delta$, и рациональный сертификат следствия~\ref{cor:poly}. Его стоимость~--- верхняя граница.
    \item Иначе построить новое отсечение и вернуться к шагу 2.
\end{enumerate}

\begin{theorem}[конечность]\label{thm:finite}
    В точной арифметике алгоритм останавливается за конечное число итераций. Итоговые числа~--- нижняя граница $M(\mathcal C)\le\rho_c(P,y)$ и верхняя граница $c(P,Q)$ при ценах $Q$ с точным сертификатом слабой разрешимости.
\end{theorem}

\begin{proof}
    Новое отсечение нарушено в $Q^*$, а все прежние в $Q^*$ выполнены. Поэтому отсечения не повторяются. Пар $(t,R)$ с $\abs R\le2$ конечное число, наборов пар тоже.
\end{proof}

\begin{remark}
    Полиномиальна только процедура разделения (теорема~\ref{thm:lp}). Мастер-задача в худшем случае экспоненциальна. Сложность вычисления $\rho_c$ при $d=2$ остаётся открытой. При $d$, входящем в данные, она NP-трудна (теорема о NP-трудности в \path{latex/perturbation.tex}). В вычислениях мастер принимает блокировку с относительным допуском $\tau=10^{-8}$, а задачи линейного программирования решаются с допуском $10^{-10}$. Тогда при $Q^*$ почти заблокированные пары выглядят строго покрывающими, и разделение повторило бы старое отсечение. Поэтому разделение делается в слегка сдвинутой точке $Q^*_\delta$, $\delta\in\set{10^{-7},10^{-6},10^{-5}}$. Корректность отсечений от точки не зависит (лемма~\ref{lem:cut}).
\end{remark}

\begin{remark}
    Метод не использует внешних решателей нелинейных задач. Нужны только задачи линейного программирования (HiGHS) и небольшие задачи с булевыми переменными для построения отсечений. Требование отрезка нужно, когда оптимум требует свидетеля с двумя ненулевыми координатами. Так в треугольнике заметок: средняя точка садится на отрезок между крайними.
\end{remark}

\section{Вычисления}\label{sec:exp}

Код лежит в \path{poly-hardness/}: \path{planar.py} (граф цепочек, задача о потоке, задача о столбце, точные сертификаты), \path{perturb2d.py} (алгоритм раздела~\ref{sec:pert}), \path{data.py}, \path{run_bw.py}, \path{exp_compare_scip.py}, тесты \path{tests/}.

\begin{itemize}
    \item \path{tests/test_planar.py}: на $400$ случайных наборах ($T\le7$, решётки с совпадениями) семейство путей графа совпадает с $\Spm(P)$, найденным по определению (пересечение отрезков свидетелей в рациональной арифметике), в обоих режимах. На $600$ парах $(P,y)$ ответ задачи о потоке совпадает с перебором конуса, задача о столбце~--- с перебором, сертификаты разрешимости проверяются точно, векторы Фаркаша разделяют.
    \item Треугольник заметок $p=(4,1),(3,3),(1,4)$, $y=(1,5,3)$: $\rho_\pi=2/5$ (известное значение), $\rho_{\abs{\log}}=2\log\frac65=0{,}36464$, $\rho_{\abs\varepsilon}=1/3$. Во всех трёх случаях оптимум требует свидетеля с двумя ненулевыми координатами.
    \item \path{exp_compare_scip.py}: на случайных наборах с $T\le6$ интервалы $[\text{нижняя},\text{верхняя}]$ алгоритма для $\pi$ пересекаются с результатами прежнего метода отсечений с точным мастером в SCIP (\path{code/perturbation/cuts.py}). RESULTS-SCIP
\end{itemize}

\paragraph{Berndt--Wood.} Шесть пар факторов из $K,L,E,M$, $T=25$, $y$~--- выпуск. Нижняя граница~--- значение мастера, верхняя проверена в рациональной арифметике.

RESULTS-BW

\section{Проверка в Lean}

Файлы \path{poly-hardness/lean/PolyHardness/FlowLP.lean} (импортирует \path{Complexity}) и \path{Blocking.lean} (импортирует \path{Algorithms.Cuts}). Проверка: \texttt{lake env lean poly-hardness/lean/PolyHardness/FlowLP.lean} и то же для \path{Blocking.lean} из корня проекта. Нет \lean{sorry}, аксиомы только стандартные (\lean{propext}, \lean{Classical.choice}, \lean{Quot.sound}).
\begin{center}
\begin{tabular}{@{}lll@{}}
\toprule
Утверждение & Файл & Lean \\
\midrule
Теорема~\ref{thm:chain}, лемма~\ref{lem:local} & \path{complexity/lean} & \lean{isReachableB_iff_chain}, \lean{inHull_iff_of_piece} и др. \\
Определение~\ref{def:graph} & \path{FlowLP} & \lean{Node}, \lean{Arc}, \lean{Out}, \lean{Bnd} \\
Ацикличность & \path{FlowLP} & \lean{exists_rank} \\
Теорема~\ref{thm:family} & \path{FlowLP} & \lean{isReachableB_iff_path} \\
Лемма~\ref{lem:disjoint} & \path{FlowLP} & \lean{path_ind_eq} \\
Лемма~\ref{lem:decomp} & \path{FlowLP} & \lean{flow_decomposition} \\
Лемма~\ref{lem:cube} & \path{FlowLP} & \lean{cube_mem_span} \\
Теорема~\ref{thm:lp} & \path{FlowLP} & \lean{mem_indicatorCone_iff_flow} \\
Утверждение~\ref{prop:min} & \path{Blocking} & \lean{exists_min_cost} \\
Лемма~\ref{lem:helly} & \path{Blocking} & \lean{isBlocked_iff_pairs} \\
Лемма~\ref{lem:pair} & \path{Blocking} & \lean{isBlocked_pair_iff}, \lean{Mixed} \\
Лемма~\ref{lem:mono} & \path{Blocking} & \lean{IsBlocked.mono} \\
Следствие~\ref{cor:cover} & \path{Blocking} & \lean{exists_pair_cover} \\
Лемма~\ref{lem:cut} & \path{Algorithms/Cuts} & \lean{exists_isBlocked_of_cut} \\
\bottomrule
\end{tabular}
\end{center}
Не формализованы: полиномиальность как утверждение о сложности (в Mathlib нет классов сложности), следствия~\ref{cor:col} и~\ref{cor:farkas} (прямые следствия теорем~\ref{thm:family} и~\ref{thm:lp} и двойственности), корректность метода ветвей и границ и конечность алгоритма (утверждение~\ref{prop:bb}, теорема~\ref{thm:finite}).

\begin{thebibliography}{99}
\bibitem{Complexity} Сложность проверки разрешимости и слабой разрешимости при фиксированных ценах, \path{complexity/complexity.tex}, 2026.
\bibitem{Pert} Минимальное возмущение цен, \path{latex/perturbation.tex}, раздел «Алгоритмическая постановка».
\end{thebibliography}

\end{document}
